\documentclass[aspectratio=169,12pt]{beamer}

\usepackage{fontspec}
\usepackage{xeCJK}
\setCJKmainfont{微軟正黑體}
\setmainfont{Times New Roman}
\XeTeXlinebreaklocale "zh"
\XeTeXlinebreakskip = 0pt plus 1pt

\usepackage{amsmath, amssymb, bm}
\usepackage{booktabs}
\usepackage{array}
\usepackage{mathtools}

\usetheme{Madrid}
\usecolortheme{default}

\title{Topic 6: Estimation}
\author{Huang, Cheng-Yang}
\institute{Statistics (III) Midterm}
\date{2026/04/21}

\begin{document}

%================================================
\begin{frame}
    \titlepage
\end{frame}

%================================================
\begin{frame}{Core Questions} % 核心問題
    \vspace{-2cm}
    \begin{itemize}
        \item What makes an estimator a ``good'' estimator?
        \item When population parameters are unknown, how do we estimate them from samples? 
        \item Why does the quality of an estimator directly affect SPC?
    \end{itemize}
\end{frame}

%================================================
\begin{frame}{Why Is Estimation Important in SPC?}
    In SPC, the in-control state of a process is usually described by baseline parameters such as $\mu_0,\, \sigma_0$

    \vspace{0.8em}
    In practice, they are usually \textbf{unknown} and must be estimated from Phase I data

    \vspace{0.8em}
    \textbf{Key points:}
        \begin{itemize}
            \item If $\hat{\mu}_0$ is estimated incorrectly, the center line will shift
            \item If $\hat{\sigma}_0$ is estimated incorrectly, the control limits will become too wide or too narrow
            \item Incorrect limits may lead to: 
            \begin{itemize}
                \item Too many false alarms
                \item Delayed detection of abnormal changes
            \end{itemize}
        \end{itemize}

    Therefore, SPC is not only about drawing charts. The first question is:

    \uncover<2->{
    {\centering
    \Huge\textbf{How good is the estimator?}\par}
    }
\end{frame}

%================================================
\begin{frame}{Estimators and Sampling Distributions}
    Suppose $$X_1, \dots,X_n \overset{iid}{\sim} F_\theta$$ where $\theta$ is an unknown parameter

    \vspace{0.5em}
    An \textbf{estimator}is a function of the sample:
    $$\hat{\theta}=T(X_1, \dots,X_n)$$ 
    Since $X_1, \dots, X_n$ are random, $\hat{\theta}$ is also a random variable

    Therefore, we should not only ask “what value do we get?”, but also ask:
        \begin{itemize}
            \item What is $E(\hat{\theta})$?
            \item What is $Var(\hat{\theta})$?
            \item Does $\hat{\theta}$ converge to $\theta$ as $n \to \infty$?
        \end{itemize}
\end{frame}

%================================================
\begin{frame}{Bias: Does the Estimator Deviate from the True Value?}
    \textbf{Definition: }
    $$\mathrm{Bias}(\hat{\theta}, \theta)=E(\hat{\theta})-\theta$$ % 長期平均誤差
    If
    $$E(\hat{\theta})=\theta,\quad \forall \theta$$ % 若長期平均誤差不偏, 則是unbiased
    then $\hat{\theta}$ is called an \textbf{unbiased estimator}

    \vspace{0.8em}
    \textbf{Interpretation: } Bias measures whether the estimator tends to overestimate or underestimate the true value on average
    
    \vspace{0.8em}
    \textbf{Example:}
    Suppose $E(X_i)=\mu$ and $\bar X=\frac1n\sum_{i=1}^n X_i.$ Then $E(\bar X)=\mu$
    
    \vspace{0.8em}
    Therefore, $\bar X$ is an unbiased estimator of $\mu$
\end{frame}

%================================================
\begin{frame}{MSE: Unbiasedness Is Not Enough}
    \vspace{-1em}
    \textbf{Mean Squared Error} is defined as
    \vspace{-0.8em}
    $$ MSE(\hat{\theta}, \theta)=E\big[(\hat{\theta}-\theta)^2\big]$$
    It can be \textbf{decomposed} into
    $$\boxed{ MSE(\hat{\theta}, \theta)=Var(\hat{\theta})+\mathrm{Bias}^2(\hat{\theta}, \theta)}$$
    \textbf{Key message: }
    \begin{itemize}
        \item An estimator can be unbiased but still have very large variance
        \item A slightly biased estimator may have a much smaller MSE
    \end{itemize}
\end{frame}

%================================================
\begin{frame}{Example}
    Goal: Estimate $\mu$, assuming\,$X_1, \dots,X_n \sim N(\mu, \sigma^2)$

    \vspace{0.5em}
    Consider two estimators: $\quad\hat{\mu}_1=\bar X, \quad\hat{\mu}_2=0.8\bar X$

    \vspace{0.5em}
    \textbf{First estimator:} $E(\hat{\mu}_1)=\mu, \quad Var(\hat{\mu}_1)=\frac{\sigma^2}{n}, \quad\text{then\quad} MSE(\hat{\mu}_1)=\frac{\sigma^2}{n}$

    \vspace{0.5em}
    \textbf{Second estimator:}
    \vspace{-0.8em}
    \begin{align*}
    &E(\hat{\mu}_2)=0.8\mu, \quad Var(\hat{\mu}_2)=0.8^2 Var(\bar X)=0.64\frac{\sigma^2}{n}\\
    &\mathrm{Bias}(\hat{\mu}_2, \mu)=0.8\mu-\mu=-0.2\mu, \\
    &\text{Then},\,  MSE(\hat{\mu}_2)=0.64\frac{\sigma^2}{n}+(0.2\mu)^2
    \end{align*}

    \textbf{Conclusion:}
    A biased estimator may still have a smaller MSE if its variance is less
\end{frame}

%================================================
\begin{frame}{Consistency: Does the Estimator Approach the True Value as the Sample Size Increases?}
    \textbf{Definition:}
    $$\hat{\theta}_n \xrightarrow{P} \theta$$
    then $\hat{\theta}_n$ is called a \textbf{consistent estimator} of $\theta$

    \vspace{0.8em}
    That is
    \vspace{-1.2em}
    $$\lim_{n\to\infty} P(|\hat{\theta}_n-\theta|>\varepsilon)=0,\quad \forall \varepsilon>0$$
    \textbf{Interpretation:}
    \begin{itemize}
        \item The estimator may have a large error when the sample size is small
        \item The estimator should become increasingly stable and closer to the true value when the sample size increases 
    \end{itemize}

    \vspace{0.8em}
    \textbf{A Common Sufficient Condition:}\quad $E(\hat{\theta}_n)\to \theta\quad\text{and}\quad Var(\hat{\theta}_n)\to 0$
\end{frame}

%================================================
\begin{frame}{Consistency of the Sample Mean}
Suppose $$X_1,\dots,X_n \overset{iid}{\sim} (\mu,\sigma^2)$$
and consider the sample mean $$\bar{X}=\frac1n\sum_{i=1}^nX_i$$
We have $$E(\bar{X})=\mu,\quad  Var(\bar{X})=\frac{\sigma^2}{n}\to0$$
Then, by Chebyshev's inequality,
$$P(|\bar{X}-\mu|>\varepsilon)\le\frac{\sigma^2/n}{\varepsilon^2}=\frac{\sigma^2}{n\varepsilon^2}\to0,\quad \forall \varepsilon>0$$
Therefore $$\bar{X}\xrightarrow{P}\mu$$
\end{frame}

\begin{frame}{Consistency of the Sample Variance}
Consider the sample variance

\vspace{-1em}
$$S^2=\frac{1}{n-1}\sum_{i=1}^n (X_i-\bar X)^2$$

\vspace{-1em}
We have

\vspace{-1em}
$$E(S^2)=\sigma^2,\quad\mathrm{Var}(S^2)\to 0$$
Then, by Chebyshev's inequality,

\vspace{-1em}
$$
P(|S^2-\sigma^2|>\varepsilon)
\le
\frac{E[(S^2-\sigma^2)^2]}{\varepsilon^2}
=
\frac{MSE(S^2)}{\varepsilon^2}
$$

Since \(E(S^2)=\sigma^2\), \(S^2\) is unbiased, so

\vspace{-1em}
$$
P(|S^2-\sigma^2|>\varepsilon)
\le
\frac{\mathrm{Var}(S^2)}{\varepsilon^2}
\to 0,
\qquad \forall \varepsilon>0
$$

Therefore
$$S^2 \xrightarrow{P} \sigma^2$$
\end{frame}

%================================================
\begin{frame}{Example}
Suppose
\vspace{-1.2em}
$$X_1, \dots,X_n \overset{iid}{\sim}\mathrm{Bernoulli}(p)$$
where $P(X_i=1)=p, \quad P(X_i=0)=1-p, \quad E(X)=p, \quad  Var(X)=p(1-p)$

\vspace{0.8em}
Since $\hat p=\bar X=\frac1n\sum_{i=1}^n X_i\Rightarrow E(\hat p)=E(\frac1n\sum_{i=1}^n X_i)=\frac1n*n*p=p$

\vspace{0.8em}
$\qquad\qquad\qquad\qquad\qquad\quad Var(\hat{p})=\frac{1}{n^2}*n*p(1-p)=\frac{p(1-p)}{n}$

\vspace{1.2em}
Therefore
\vspace{-1.4em}
$$E(\hat p)\to p\,\,\,\text{and}\,\,\, Var(\hat p)=\frac{p(1-p)}{n}\to 0\quad \text{as}\quad n\to\infty$$
So $$\hat{p}\overset{P}{\longrightarrow}p\Rightarrow\hat{p}\,\,\, \text{is consistent}$$ 
\end{frame}

%================================================
\begin{frame}{From Estimation to SPC: The Core of Phase I}
    In Phase I, we first establish an in-control baseline

    \vspace{0.5em}
    For an $\bar{X}$-chart, if $\mu_0$ and $\sigma_0$ are known, the control limits are
    $$
    UCL=\mu_0+z_{1-\alpha/2}\frac{\sigma_0}{\sqrt{n}},
    \quad
    LCL=\mu_0-z_{1-\alpha/2}\frac{\sigma_0}{\sqrt{n}}
    $$
    Usually, we replace them with estimators $\hat{\mu}_0$ and $\hat{\sigma}_0$
    
    \vspace{0.5em}
    Thus, the actual control limits become
    $$
    \widehat{UCL}=\hat{\mu}_0+z_{1-\alpha/2}\frac{\hat{\sigma}_0}{\sqrt{n}},
    \quad
    \widehat{LCL}=\hat{\mu}_0-z_{1-\alpha/2}\frac{\hat{\sigma}_0}{\sqrt{n}}
    $$
    \textbf{Potential issues:}
    \begin{itemize}
        \item If $\hat{\mu}_0$ is biased, the center line shifts.
        \item If $\hat{\sigma}_0$ has large variance, the control limits become unstable.
        \item If the estimators are not consistent, we cannot establish a reliable baseline.
    \end{itemize}
\end{frame}

%================================================
\begin{frame}{How Estimator Affects Monitoring: Example of a p-Chart}
    Suppose each subgroup has sample size $n$, and the defect rate is $p_0$
    Then $$\hat p=\frac{X}{n}, \quad X\sim\mathrm{Binomial}(n, p_0)$$
    In Phase I, we usually use $\bar{p}=\frac1m\sum^m_{i=1}\hat{p_i}$ to estimate $p_0$ 

    The control limits of the p-chart are often written as

    $$
    UCL=\bar p + z\sqrt{\frac{\bar p(1-\bar p)}{n}},
    \quad
    LCL=\bar p - z\sqrt{\frac{\bar p(1-\bar p)}{n}}
    $$
\end{frame}

%================================================
\begin{frame}{Conclusion}
    \vspace{0.2cm}
    \textbf{(1) A good estimator should not be judged only by unbiasedness}
    \begin{itemize}
        \item $ MSE= Var+\mathrm{Bias}^2$
    \end{itemize}

    \vspace{0.8em}
    \textbf{(2) Consistency is a requirement for long-run reliability}
    \begin{itemize}
        \item $\hat{\theta}_n \xrightarrow{P} \theta$
    \end{itemize}

    \vspace{0.8em}
    \textbf{(3) In SPC, estimation is the foundation}
    \begin{itemize}
        \item Phase I requires estimation of $\mu_0$ and $\sigma_0$ in order to determine control limits
        \item Control limits determine false alarm rates and detection power
    \end{itemize}
    \vspace{0.8em}
    \uncover<2->{$$\boxed{\begin{array}{c}\text{If the baseline estimation is unreliable,} \\\text{then the entire SPC monitoring system becomes unreliable.}\end{array}}$$}
\end{frame}

\begin{frame}
    \centering\Huge{\textbf{Thank You}}
\end{frame}

\end{document}